\documentclass[11pt]{article}
\usepackage[margin=1in]{geometry}
\usepackage{booktabs}
\usepackage{amsmath}
\usepackage{array}
\usepackage{hyperref}

\title{ECE-GY 9143 Lab 4 Part A Report\\Distributed Deep Learning}
\author{Haojingtong Tong}
\date{\today}

\begin{document}
\maketitle

\section{Experiment Setup}
\textbf{Hardware:} NVIDIA A100-SXM4-80GB (up to 4 GPUs).\\
\textbf{Model/Data:} ResNet-18 on CIFAR-10.\\
\textbf{Optimizer:} SGD, learning rate $0.1$, momentum $0.9$, weight decay $5\times 10^{-4}$, num\_workers $=2$.\\
\textbf{Protocol:} Each timing experiment runs 2 epochs and reports the 2nd epoch.

\section{Q1: Computational Efficiency vs Batch Size}
Table~\ref{tab:q1} reports the 2nd-epoch training time (excluding data loading) for single-GPU runs.

\begin{table}[h]
\centering
\caption{Q1 single-GPU training time (2nd epoch, excluding data I/O)}
\label{tab:q1}
\begin{tabular}{cc}
\toprule
Batch size per GPU & Train time (s) \\
\midrule
32   & 12.543 \\
128  & 5.348 \\
512  & 4.350 \\
2048 & 4.152 \\
8192 & 4.135 \\
\bottomrule
\end{tabular}
\end{table}

At batch size 32768, the run hit CUDA OOM. Therefore, the largest feasible single-GPU batch size is \textbf{8192}.  
As batch size increases, runtime decreases sharply at first (better GPU utilization), then plateaus around $\sim 4.1$s.

\section{Q2: Speedup Measurement}
We measured with 2 and 4 GPUs for batch sizes 32/128/512 (per GPU).  
Speedup is computed as:
\[
\text{Speedup}(N)=\frac{T_{1\text{GPU}}}{T_{N\text{GPU}}}
\]
using 2nd-epoch \texttt{total\_time}.

\begin{table}[h]
\centering
\caption{Table 1: Training time and speedup}
\label{tab:q2}
\begin{tabular}{c|cc|cc|cc}
\toprule
 & \multicolumn{2}{c|}{Batch 32 / GPU} & \multicolumn{2}{c|}{Batch 128 / GPU} & \multicolumn{2}{c}{Batch 512 / GPU} \\
GPUs & Time (s) & Speedup & Time (s) & Speedup & Time (s) & Speedup \\
\midrule
1-GPU & 13.340 & 1.000 & 5.995 & 1.000 & 5.808 & 1.000 \\
2-GPU & 8.690  & 1.535 & 3.253 & 1.843 & 3.095 & 1.877 \\
4-GPU & 4.540  & 2.938 & 1.775 & 3.377 & 1.644 & 3.533 \\
\bottomrule
\end{tabular}
\end{table}

\paragraph{Scaling type discussion.}
This experiment is \textbf{weak scaling}: per-GPU batch size is fixed while effective global batch size increases with GPU count.  
If we measured strong scaling (fixed global batch size), speedup would likely be worse because each GPU would process smaller local batches and communication overhead would become relatively larger.  
A supporting data point from Table~\ref{tab:q2}: speedup is lower at smaller per-GPU batch (e.g., 4-GPU speedup is 2.938 at batch 32 vs 3.533 at batch 512).

\section{Q3: Computation vs Communication}
\subsection*{Q3.1 Compute and Communication Time}
To decompose DDP train time, we use:
\[
T_{\text{compute}}(N,b)\approx \frac{T^{\text{Q1}}_{\text{train}}(b)}{N},\qquad
T_{\text{comm}}(N,b)=T^{\text{Q3}}_{\text{train}}(N,b)-T_{\text{compute}}(N,b)
\]
where $N$ is number of GPUs and $b$ is batch size per GPU.

\begin{table}[h]
\centering
\caption{Table 2: Compute and communication time (seconds)}
\label{tab:q3_1}
\begin{tabular}{c|cc|cc|cc}
\toprule
 & \multicolumn{2}{c|}{Batch 32 / GPU} & \multicolumn{2}{c|}{Batch 128 / GPU} & \multicolumn{2}{c}{Batch 512 / GPU} \\
GPUs & Compute & Comm & Compute & Comm & Compute & Comm \\
\midrule
2-GPU & 6.272 & 1.870 & 2.674 & 0.376 & 2.175 & 0.093 \\
4-GPU & 3.136 & 1.270 & 1.337 & 0.260 & 1.088 & 0.348 \\
\bottomrule
\end{tabular}
\end{table}

\subsection*{Q3.2 Communication Bandwidth Utilization}
For ring all-reduce in \cite{patarasuk2009}, per-iteration communication volume is:
\[
V_{\text{AR}} = 2\frac{N-1}{N}S
\]
where $S$ is model size in bytes. Here, $S=44{,}695{,}848$ bytes ($\approx 0.0416$ GB).  
Per-epoch effective communication bandwidth is computed as:
\[
\text{BW}_{\text{eff}}=\frac{V_{\text{AR}}\times \#\text{iters per epoch}}{T_{\text{comm,epoch}}}
\]

\begin{table}[h]
\centering
\caption{Table 3: Effective communication bandwidth utilization (GB/s)}
\label{tab:q3_2}
\begin{tabular}{c|c|c|c}
\toprule
GPUs & Batch 32 / GPU & Batch 128 / GPU & Batch 512 / GPU \\
\midrule
2-GPU & 17.412 & 21.699 & 21.932 \\
4-GPU & 19.220 & 23.535 & 4.492 \\
\bottomrule
\end{tabular}
\end{table}

\section{Q4: Large Batch Training}
\subsection*{Q4.1 Accuracy with large batch}
From Q1, the largest feasible batch size per GPU is 8192.  
Running 4 GPUs for 5 epochs gives:
\[
\text{Effective batch size}=4\times8192=32768
\]
\[
\text{Epoch 5 avg loss}=3.1467,\quad
\text{Epoch 5 avg accuracy}=18.31\%
\]

For comparison, the batch-128 single-GPU run (Q1 measured epoch) had:
\[
\text{avg loss}=1.3840,\quad \text{avg top-1}=49.49\%
\]
The large-batch setting significantly degrades accuracy without additional large-batch training techniques.

\subsection*{Q4.2 Two remedies to improve large-batch accuracy}
Based on \cite{goyal2017}:
\begin{enumerate}
    \item \textbf{Linear learning-rate scaling with warmup:} scale learning rate proportionally to batch-size increase and use gradual warmup in early epochs.
    \item \textbf{Careful BatchNorm handling in distributed training:} maintain consistent BN behavior/statistics across workers to stabilize optimization.
\end{enumerate}

\section{Q5: What is passed on network?}
\textbf{Not only gradients.}  
In DDP, gradient all-reduce is the main per-iteration communication, but model consistency also involves synchronization operations such as initial parameter broadcast and buffer synchronization (e.g., BatchNorm running statistics, depending on configuration).

\section{Q6: What if we only communicate gradients?}
For 4 GPUs with batch size 512 (or larger), communicating only gradients is \textbf{not sufficient} to guarantee high accuracy.  
Large-batch training quality also depends on optimization strategy (e.g., LR scaling and warmup) and proper handling of normalization/statistics as discussed in \cite{goyal2017}. Synchronization alone does not solve large-batch optimization instability.

\begin{thebibliography}{9}
\bibitem{dist_overview}
PyTorch, \emph{Distributed Overview}.\\
\url{https://pytorch.org/tutorials/beginner/dist_overview.html}

\bibitem{ddp_doc}
PyTorch, \emph{torch.nn.parallel.DistributedDataParallel}.\\
\url{https://docs.pytorch.org/docs/stable/generated/torch.nn.parallel.DistributedDataParallel.html}

\bibitem{patarasuk2009}
P. Patarasuk and X. Yuan, ``Bandwidth optimal all-reduce algorithms for clusters of workstations,'' \emph{J. Parallel Distrib. Comput.}, vol. 69, pp. 117--124, 2009.

\bibitem{goyal2017}
P. Goyal et al., ``Accurate, Large Minibatch SGD: Training ImageNet in 1 Hour,'' \emph{arXiv:1706.02677}, 2017.
\end{thebibliography}

\end{document}
